% Part: model-theory
% Chapter: interpolation
% Section: definability

\documentclass[../../../include/open-logic-section]{subfiles}

\begin{document}

\olfileid[pt]{mod}{int}{def}

\olsection{O Teorema da Definibilidade}

Uma aplicação importante do teorema da interpolação é o teorema da
definibilidade de Beth. Para definir uma relação $n$-ária~$R$, podemos
fornecer !!a{formula}~$!C$ com $n$ !!{variable}s livres que não
envolva~$R$. Isso seria uma definição \emph{explícita} de~$R$ em termos
de~$!C$. Podemos também dizer que uma teoria~$\Sigma(P)$ em uma
!!{language} que contém o !!{predicate} $n$-ário~$P$ define
explicitamente~$P$ se contém (ou ao menos implica) uma definição
explícita formalizada, O.s.\@,
\[
\Sigma(P) \Entails \lforall[x_1][\dots
  \lforall[x_n][(\Atom{P}{x_1,\dots, x_n} \liff !C(x_1, \dots,
    x_n))]].
\]
Mas uma definição explícita é apenas uma maneira de definir---no
sentido de determinar por completo---uma relação. Uma teoria também
pode ser tal que, em qualquer modelo, a interpretação de~$P$ seja
fixada pela interpretação do restante da !!{language}. O teorema da
definibilidade afirma que, sempre que uma teoria fixa a interpretação
de~$P$ dessa maneira---sempre que ela \emph{define implicitamente}
~$P$---ela também o define explicitamente.

\begin{defn}
Suponha que $\Lang{L}$ seja !!a{language} que não contém o
!!{predicate}~$P$. Um conjunto $\Sigma(P)$ de !!{sentence}s de
$\Lang{L} \cup \{P\}$ \emph{define explicitamente}~$P$ se, e somente
se, existe !!a{formula}~$!C(x_1, \dots, x_n)$ de $\Lang{L}$ tal que
\[
\Sigma(P) \Entails \lforall[x_1][\dots
  \lforall[x_n][(\Atom{P}{x_1,\dots, x_n} \liff !C(x_1, \dots,
    x_n))]].
\]
\end{defn}

\begin{defn}
Suponha que $\Lang{L}$ seja !!a{language} que não contém os
!!{predicate}s~$P$ e~$P'$. Um conjunto $\Sigma(P)$ de !!{sentence}s de
$\Lang{L} \cup \{P\}$ \emph{define implicitamente} $P$ se, e somente
se,
\[
\Sigma(P) \cup \Sigma(P') \Entails \lforall[x_1][\dots
  \lforall[x_n][(\Atom{P}{x_1,\dots, x_n} \liff \Atom{P'}{x_1,\dots,
      x_n})]],
\]
onde $\Sigma(P')$ resulta da substituição uniforme de $P$ por $P'$ em
$\Sigma(P)$.
\end{defn}

Em outras palavras, para qualquer modelo $\Struct{M}$ e quaisquer
$R, R' \subseteq \Domain{M}^n$, se tanto
$\Expan{M}{R} \Entails \Sigma(P)$ quanto
$\Expan{M}{R'} \Entails \Sigma(P')$, então $R=R'$; aqui,
$\Expan{M}{R}$ é a !!{structure}~$\Struct{M'}$ para a expansão de
$\Lang{L}$ a $\Lang{L} \cup \{P\}$ tal que $\Assign{P}{M'} = R$, e de
modo análogo para $\Expan{M}{R'}$.

\begin{thm}[Teorema da Definibilidade de Beth] Um conjunto $\Sigma(P)$
  de $\Lang{L} \cup\{P\}$-!!{formula}s define implicitamente $P$ se, e
  somente se, $\Sigma(P)$ define explicitamente $P$.
\end{thm}

\begin{proof}
Se $\Sigma(P)$ define explicitamente $P$, então valem ambas as relações
\begin{align*}
  \Sigma(P) & \Entails & \lforall[x_1][\dots \lforall[x_n]
    [(\Atom{P}{x_1,\dots, x_n} \liff !C(x_1,\dots,x_n))]]\\
  \Sigma(P') & \Entails & \lforall[x_1][\dots \lforall[x_n]
    [(\Atom{P'}{x_1,\dots, x_n} \liff !C(x_1,\dots,x_n))]]
\end{align*}
e a conclusão se segue. Para a recíproca, suponha que $\Sigma(P)$
defina implicitamente $P$. Primeiro, acrescentamos as !!{constant}s
$c_1$, \dots,~$c_n$ a $\Lang{L}$. Então
\[
\Sigma(P) \cup \Sigma(P') \Entails
\Atom{P}{c_1, \dots, c_n} \to  \Atom{P'}{c_1, \dots, c_n}.
\]
Pela compacidade, existem conjuntos finitos
$\Delta_0 \subseteq \Sigma(P)$ e $\Delta_1 \subseteq \Sigma(P')$ tais
que
\[
\Delta_0 \cup \Delta_1 \Entails
\Atom{P}{c_1, \dots, c_n} \to \Atom{P'}{c_1, \dots, c_n}.
\]
Seja $!D(P)$ a conjunção de todas as !!{sentence}s $!A(P)$ tais que
$!A(P) \in \Delta_0$ ou $!A(P') \in \Delta_1$, e seja $!D(P')$ a
conjunção de todas as !!{sentence}s $!A(P')$ tais que
$!A(P) \in \Delta_0$ ou $!A(P') \in \Delta_1$. Então
$!D(P) \land !D(P') \Entails
\Atom{P}{c_1, \dots, c_n} \to P'c_1\dots c_n$. Podemos reorganizar
essa relação de modo que cada !!{predicate} ocorra de um lado de
$\Entails$:
\[
!D(P) \land \Atom{P}{c_1, \dots, c_n} \Entails
!D(P') \to \Atom{P'}{c_1, \dots, c_n}.
\]
Pelo Teorema da Interpolação de Craig, existe !!a{sentence}
$!C(c_1,\dots, c_n)$ que não contém $P$ nem $P'$ e é tal que:
\begin{align*}
  !D(P) \land \Atom{P}{c_1, \dots, c_n} & \Entails !C(c_1,\dots, c_n); \\
  !C(c_1,\dots, c_n) & \Entails !D(P') \to \Atom{P'}{c_1, \dots, c_n}.
\end{align*}
Da primeira dessas duas relações de consequência, obtemos:
$!D(P) \Entails
\Atom{P}{c_1,\dots, c_n} \lif !C(c_1,\dots, c_n)$. E, da segunda,
como um modelo de $\Lang{L} \cup \{P\}$, $\Expan{M}{R}
\Entails !A(P)$ se, e somente se, o modelo correspondente de
$\Lang{L} \cup \{P'\}$, $\Expan{M}{R} \models !A(P')$, obtemos
$!C(c_1,\dots, c_n) \Entails
!D(P) \lif \Atom{P}{c_1,\dots, c_n}$, donde:
\[
!D(P) \Entails !C(c_1,\dots,c_n) \to \Atom{P}{c_1,\dots, c_n}.
\]
Reunindo as duas relações,
$!D(P) \Entails \Atom{P}{c_1,\dots, c_n}
\liff !C(c_1, \dots, c_n)$, e, por monotonicidade e generalização,
também
\[
\Sigma(P) \Entails
\lforall[x_1][\dots\lforall[x_n][(\Atom{P}{x_1,\dots, x_n} \liff
    !C(x_1,\dots, x_n))]].
\]
\end{proof}

\end{document}

